%=====================================================================
\chapter{Research Methods and Experimental Design}\label{ch:research}
%=====================================================================
\locator{6}{Part VI --- Professional Practice}

\begin{outcomes}
By the end of this chapter you will be able to:
\begin{tight}
  \item \textbf{Convert} a broad topic into a researchable question with defined scope.
  \item \textbf{Select} a research design appropriate to the question type.
  \item \textbf{Distinguish} the forms of validity and reliability, and state a threat to each.
  \item \textbf{Structure} a research proposal and a final report.
  \item \textbf{Apply} reproducibility practices to your own work.
\end{tight}
\end{outcomes}

\begin{prereq}
\chref{inference} (hypothesis testing, power, multiple comparisons) and
\chref{causality} (designs). This chapter is the framework those techniques sit in,
and is what your final-year project will be assessed against.
\end{prereq}

\begin{keyterms}
research question $\bullet$ literature review $\bullet$ hypothesis $\bullet$
independent and dependent variable $\bullet$ operationalisation $\bullet$
internal / external / construct / conclusion validity $\bullet$ reliability
$\bullet$ triangulation $\bullet$ pre-registration $\bullet$ reproducibility
$\bullet$ replication
\end{keyterms}

\section{From Topic to Question}

\begin{alertbox}[A Topic Is Not a Question]
``Machine learning in education'' is a topic. You cannot answer it, cannot tell
whether you have finished, and cannot fail. A researchable question has a
population, a variable, an outcome and a scope --- and it is possible to be wrong
about it.
\end{alertbox}

\begin{worked}[Narrowing a Topic, Four Times]
\textbf{Round 0 --- topic.} ``AI in education.'' Unanswerable.

\smallskip
\textbf{Round 1 --- add a phenomenon.} ``Does AI help students?'' Which AI? Which
students? Help how?

\smallskip
\textbf{Round 2 --- add a population and an outcome.} ``Do first-year BSIT students
who use an AI tutor get better grades?'' Better than whom? Which grades? Over what
period?

\smallskip
\textbf{Round 3 --- operationalise everything.} ``Among first-year BSIT students at
one institution in Fall 2026, does access to a retrieval-grounded AI tutor increase
the end-of-semester programming module mark, compared with the existing
drop-in-help provision?''

\smallskip
\textbf{Round 4 --- check it is answerable.}
\begin{tight}
  \item \emph{Population:} first-year BSIT, one institution, one semester --- so
        external validity is limited, and you will say so.
  \item \emph{Independent variable:} access to the tutor (yes/no), assigned how?
        Randomly, or you have a confounding problem (\chref{causality}).
  \item \emph{Dependent variable:} final module mark out of 100 --- measurable,
        recorded already.
  \item \emph{Comparison:} existing provision, not nothing. Comparing against
        nothing measures ``having help'' rather than ``having \emph{this} help''.
  \item \emph{Feasible?} About 630 students needed for a 10-point effect at 80\%
        power (\chref{causality}). The cohort is 340. \textbf{So state honestly
        that the study is powered to detect roughly 14 points, not 10} --- or
        pre-register a two-semester design.
\end{tight}

\smallskip
\textbf{The last check is the one students skip}, and it is why so many projects
conclude ``no significant difference'' from a study that never could have found
one.
\end{worked}

\section{Choosing a Design}

\dsfig{%
\begin{tikzpicture}[font=\scriptsize,
  q/.style={rounded corners=3pt, draw=primary, fill=primary!8, text=primary!55!black,
            text width=3.0cm, align=center, font=\tiny\bfseries, minimum height=0.8cm},
  d/.style={rounded corners=3pt, draw=#1, fill=#1!8, text=#1!50!black,
            text width=3.0cm, align=center, font=\tiny, minimum height=1.15cm},
  arr/.style={-{Stealth[length=2mm]}, primary!50, line width=0.85pt}]

\node[q] (root) at (5.4,1.4) {What kind of question is it?};

\node[d=secondary] (a) at (0,-0.5)   {\textbf{``What is happening?''}\\descriptive /\\
                                      survey study};
\node[d=greenhl]   (b) at (3.6,-0.5) {\textbf{``Are these related?''}\\correlational /\\
                                      observational};
\node[d=goldhl]    (c) at (7.2,-0.5) {\textbf{``Does X cause Y?''}\\experiment or\\
                                      quasi-experiment};
\node[d=purplehl]  (e) at (10.8,-0.5){\textbf{``Why / how?''}\\qualitative:\\
                                      interviews, cases};

\foreach \n in {a,b,c,e}{\draw[arr] (root) -- (\n);}

\node[font=\tiny, text=gray!55!black, align=center, text width=3.1cm] at (0,-2.0)
     {\emph{Cannot} establish cause. Watch sampling bias (\chref{inference}).};
\node[font=\tiny, text=gray!55!black, align=center, text width=3.1cm] at (3.6,-2.0)
     {\emph{Cannot} establish cause either. Report confounders you could not control.};
\node[font=\tiny, text=gray!55!black, align=center, text width=3.1cm] at (7.2,-2.0)
     {The only design that supports a causal claim. State the assumption if not
      randomised.};
\node[font=\tiny, text=gray!55!black, align=center, text width=3.1cm] at (10.8,-2.0)
     {Generates hypotheses and explains mechanism; not generalisable by itself.};

\node[anchor=north, rounded corners=3pt, draw=tealhl, fill=tealhl!8, text=tealhl!50!black,
      text width=12.6cm, align=center, font=\scriptsize, minimum height=0.65cm] at (5.4,-3.1)
     {\textbf{MIXED METHODS:} quantitative shows \emph{that} something happens;
      qualitative explains \emph{why}. Using both is \term{triangulation}, and it is
      usually the strongest design available to a student project.};
\end{tikzpicture}}%
{Choosing a design from the question. The commonest error in student projects is running
a correlational study and writing a causal conclusion.}{research-designs}

\section{Validity and Reliability}

\rowcolors{2}{white}{lightbg}
\begin{longtable}{@{}L{3.0cm}L{4.6cm}L{6.4cm}@{}}
\toprule
\rowcolor{primary!15}
\textbf{Type} & \textbf{Asks} & \textbf{A threat, and its remedy} \\
\midrule
\endhead
Internal validity & Did X really cause Y, here? & Confounding. \emph{Remedy:} randomise, or control and state the assumption \\
External validity & Does it generalise beyond this sample? & One cohort at one institution. \emph{Remedy:} describe the context precisely; do not overclaim \\
Construct validity & Are you measuring what you claim? & ``Engagement'' measured as clicks --- a student may click without engaging. \emph{Remedy:} multiple indicators \\
Conclusion validity & Are the statistical inferences sound? & Underpowered study; multiple comparisons. \emph{Remedy:} power analysis first, corrections after (\chref{inference}) \\
Reliability & Would you get the same result again? & Inconsistent coding of qualitative data. \emph{Remedy:} inter-rater agreement, published rubric \\
\bottomrule
\end{longtable}

\begin{conceptbox}[Construct Validity Is the One Data Scientists Get Wrong]
We are trained to optimise a number, so we accept whatever number is available. But
\texttt{clicks} is not engagement, \texttt{commits} is not productivity,
\texttt{time on page} is not attention, and \texttt{alerts closed} is not security.
Each is a \emph{proxy}, and the gap between the proxy and the construct is where
projects quietly go wrong --- especially once people learn they are measured by it.
Name the construct, name the proxy, and state the gap.
\end{conceptbox}

\section{Reproducibility}

\dsfig{%
\begin{tikzpicture}[font=\scriptsize,
  lvl/.style={rounded corners=3pt, draw=#1, fill=#1!8, text=#1!50!black,
              text width=3.35cm, align=center, font=\scriptsize, inner sep=5pt,
              minimum height=2.5cm},
  hd/.style={rounded corners=3pt, fill=#1, text=white, font=\scriptsize\bfseries,
             text width=3.35cm, align=center, minimum height=0.5cm}]

\node[hd=accent] at (0,0)     {NOT REPRODUCIBLE};
\node[lvl=accent] at (0,-1.65){%
  A results table in a report.\\[4pt]
  Nobody, including you in six months, can regenerate it.\\[4pt]
  \emph{Sadly, the default.}};

\node[hd=goldhl] at (3.9,0)     {REPRODUCIBLE};
\node[lvl=goldhl] at (3.9,-1.65){%
  Same data + same code $\rightarrow$ same result.\\[4pt]
  Needs: versioned code, data snapshot, pinned environment, fixed random seeds.\\[4pt]
  \emph{The minimum bar.}};

\node[hd=secondary] at (7.8,0)     {REPLICABLE};
\node[lvl=secondary] at (7.8,-1.65){%
  New data + same method $\rightarrow$ same conclusion.\\[4pt]
  Needs: a method described precisely enough for someone else to follow.\\[4pt]
  \emph{What science actually requires.}};

\node[hd=greenhl] at (11.7,0)     {ROBUST};
\node[lvl=greenhl] at (11.7,-1.65){%
  Different reasonable analyses $\rightarrow$ same conclusion.\\[4pt]
  Needs: sensitivity analysis --- try other thresholds, other models, other
  exclusions.\\[4pt]
  \emph{The strongest claim available.}};

\node[anchor=north west, align=left, text width=13.4cm, font=\scriptsize, text=primary!60!black]
     at (-1.9,-3.5)
     {\textbf{The cheapest single improvement to a student project:} fix every random
      seed, pin every library version in a container (\chref{cloudedge}), and commit the
      analysis script that produced each figure. This costs an hour and moves you a whole
      column to the right.};
\end{tikzpicture}}%
{Four levels of confidence in a result. Most student and industry work sits in the first
column; an hour of engineering moves it to the second.}{reproducibility}

\begin{definitionbox}[Pre-registration]
Recording your hypothesis, method, sample size and analysis plan \emph{before}
collecting or examining the data. It is the structural defence against $p$-hacking
(\chref{inference}): decisions made in advance cannot be adjusted to fit the result.
Even an unpublished, dated document in your repository provides most of the benefit.
\end{definitionbox}

\section{Structuring the Write-Up}

\rowcolors{2}{white}{defbg}
\begin{longtable}{@{}L{2.7cm}L{4.4cm}L{6.9cm}@{}}
\toprule
\rowcolor{primary!15}
\textbf{Section} & \textbf{Answers} & \textbf{Common failure} \\
\midrule
\endhead
Abstract & The whole study in 200 words & Written first. Write it last \\
Introduction & Why does this matter? & Background with no stated gap or question \\
Literature review & What is already known? & A list of summaries instead of a synthesis that leads to your gap \\
Methodology & What did you do, exactly? & Not enough detail for anyone to repeat it \\
Results & What did you find? & Interpretation smuggled into the results section \\
Discussion & What does it mean? & Overclaiming; ignoring limitations \\
Limitations & Where is it weak? & Omitted, or reduced to ``more data would help'' \\
Conclusion & So what? & New material appearing for the first time \\
\bottomrule
\end{longtable}

\begin{conceptbox}[The Limitations Section Is Where You Gain Credibility]
Students treat limitations as a confession that weakens the work. Examiners read it
as evidence that you understand your own study. ``The sample was one cohort at one
institution, so generalisation to other contexts is not supported; the study was
powered to detect a 14-point effect, so a smaller true effect would have been
missed'' demonstrates more competence than any result in the paper. Write it
specifically, with numbers.
\end{conceptbox}

%---------------------------------------------------------------------
\begin{fourlenses}{Research Design}
\lensLA{\emph{Question:} does a specific intervention improve a specific outcome?
\emph{Design:} randomise which students are \emph{offered} the intervention, then
analyse by offer rather than by uptake --- this keeps the comparison causal even
though attendance is voluntary. \emph{Validity threat:} construct validity, because
``engagement'' is measured by clicks. \emph{Ethical constraint:} a control group
denied a support service believed to work is hard to justify, so use a stepped-wedge
design in which every group receives it eventually (\chref{ethics}).}

\lensPM{\emph{Question:} does a practice change delivery predictability?
\emph{Design:} randomising methodology across projects is rarely feasible, so use
difference-in-differences across teams adopting at different times
(\chref{causality}). \emph{Validity threat:} internal validity --- teams that adopt
a new practice early differ systematically from those that do not.
\emph{Sample-size reality:} with 12 teams, no quantitative result will be
conclusive, so pair it with interviews and report it as a mixed-methods case study.}

\lensIOT{The easiest domain in this book for clean experiments: devices can be
randomised at essentially no cost and in large numbers, with no consent problem.
Push firmware A to a random half of a 400-device fleet and you have a genuine RCT
with $n=400$. \emph{Validity threat:} external validity --- results from one site's
equipment, duty cycle and environment may not transfer to another. Report the
operating conditions as carefully as the results.}

\lensSEC{\emph{Question:} does a detection change reduce time to detect?
\emph{Design:} A/B test rules on randomly split traffic; observational comparison is
almost worthless here because organisations that buy more tooling also suffer more
incidents (reverse causation, \chref{causality}). \emph{Validity threat:} conclusion
validity --- attacks are rare, so studies are chronically underpowered, and a
non-significant result usually means the study was too small rather than that
nothing happened. \emph{Special constraint:} publishing method detail helps
attackers, which creates a genuine tension with reproducibility that must be managed
rather than ignored.}
\end{fourlenses}

%---------------------------------------------------------------------
\begin{lab}[Making Your Own Project Reproducible]
\begin{lstlisting}[language=Python]
"""Put this at the top of every analysis script. It costs five minutes."""
import random, os, json, subprocess
from datetime import datetime
import numpy as np

SEED = 42
random.seed(SEED); np.random.seed(SEED)
os.environ["PYTHONHASHSEED"] = str(SEED)
# torch.manual_seed(SEED); torch.use_deterministic_algorithms(True)

def provenance(outfile="run_provenance.json", **params):
    """Record everything needed to regenerate this result."""
    rec = {
        "timestamp":   datetime.utcnow().isoformat(),
        "git_commit":  subprocess.check_output(
                           ["git", "rev-parse", "HEAD"]).decode().strip(),
        "git_dirty":   bool(subprocess.check_output(
                           ["git", "status", "--porcelain"]).decode().strip()),
        "seed":        SEED,
        "python":      subprocess.check_output(
                           ["python", "--version"]).decode().strip(),
        "params":      params,
    }
    if rec["git_dirty"]:
        print("WARNING: uncommitted changes -- this run is not reproducible")
    json.dump(rec, open(outfile, "w"), indent=2)
    return rec

provenance(model="logreg", threshold=0.35, features=len(FEATURES))
\end{lstlisting}

\textbf{A pre-registration you can write in ten minutes} --- commit it before you
look at the data:
\begin{lstlisting}[basicstyle=\ttfamily\tiny, frame=none, backgroundcolor=\color{codebg}]
QUESTION   Among first-year BSIT students (Fall 2026, n=340), does access to
           the AI tutor increase final programming mark vs existing provision?
H0         No difference in mean final mark.
H1         The tutor group scores higher.
DESIGN     Randomised at the point of OFFER; analysed by offer (ITT).
PRIMARY    Final module mark /100. ONE primary outcome only.
SECONDARY  Submission count; withdrawal rate. Reported as exploratory.
ANALYSIS   Two-sample t-test, alpha = 0.05, two-sided. Cohen's d reported.
POWER      n=340 gives 80% power to detect ~14 marks. Smaller effects will
           not be detectable and a null result must be reported as such.
EXCLUSIONS Students withdrawing before week 3. Defined NOW, not later.
DATE       2026-09-01, before any data was examined.
\end{lstlisting}
\end{lab}

%---------------------------------------------------------------------
\begin{checkpoint}
\begin{tightnum}
  \item Turn ``blockchain in supply chains'' into a researchable question with a
        population, variables and a scope.
  \item You measure ``student engagement'' by number of VLE clicks. Which validity
        type is threatened, and how would you strengthen it?
  \item What is the difference between a reproducible and a replicable result?
\end{tightnum}
\end{checkpoint}

%---------------------------------------------------------------------
\begin{chaptersummary}
\begin{tight}
  \item A researchable question names a population, an independent and a dependent
        variable, a comparison and a scope --- and is capable of being wrong.
  \item The design follows from the question type; only experiments and
        quasi-experiments support causal claims.
  \item Check power \emph{before} running the study, and report what effect size you
        were actually able to detect.
  \item Validity comes in internal, external, construct and conclusion forms.
        Construct validity is the one data scientists most often neglect.
  \item Reproducible $\rightarrow$ replicable $\rightarrow$ robust. Seeds, pinned
        environments and committed scripts move you up cheaply.
  \item Pre-register the hypothesis, sample size, primary outcome and exclusions
        before looking at the data.
  \item A specific, numerical limitations section adds credibility rather than
        removing it.
\end{tight}
\end{chaptersummary}

\begin{reviewq}
  \item \textbf{(MCQ)} Which design can support a causal claim? (a)~survey
        (b)~correlational study (c)~randomised experiment (d)~case study.
  \item \textbf{(MCQ)} Measuring ``productivity'' by lines of code primarily
        threatens: (a)~internal (b)~external (c)~construct (d)~conclusion validity.
  \item \textbf{(Short)} Define pre-registration and explain which specific problem
        from \chref{inference} it prevents.
  \item \textbf{(Short)} Explain intention-to-treat analysis and why it preserves a
        causal interpretation when uptake is voluntary.
  \item \textbf{(Short)} Give two reasons a non-significant result is not evidence
        that no effect exists.
  \item \textbf{(Applied)} Design a study to test whether a new code-review practice
        reduces defects. State the question, design, variables, sample size
        reasoning, one threat to each of the four validity types, and how you would
        make it reproducible.
  \item \textbf{(Applied)} A colleague reports ``students using the library score
        8 marks higher, so we should require library use''. Identify every
        methodological problem and rewrite the conclusion defensibly.
\end{reviewq}
